\documentclass[a4paper, 12pt]{article}

\usepackage{utils}

\renewcommand*{\today}{14 November 2025}

\begin{document}

\hotbox{Géometrie 1}{CM 8}{\today}

\begin{demonstration}
    \begin{enumerate}
        \item \dots
        \item Soit $f: \E \to \mathcal{F}$ une application affine. Montrons que $f^{-1}: \mathcal{F} \to \E$ existe,
        $f^{-1}$ est affine et que $\vec{f^{-1}} = \vec{f}^{-1}$.

        Soit $b \in \mathcal{F}$. L'application $(\vec{f^{-1}})_b$ associée à $f^{-1}$ et au point $b$ est définie par :
        $$
        (\vec{f^{-1}})_b (\vec{bn}) = \vec{f^{-1}(b) f^{-1}(n)}
        $$
        Soit $a \in \E$ et $m \in \E$ tels que $b = f(a)$ et $n = f(m)$. D'où $(\vec{f^{-1}})_b (\vec{bn}) = \vec{am}$.
        Or, $\vec{f}(\vec{am}) = \vec{f(a)f(m)} = \vec{bn}$ donc comme $\vec{f}$ est bijective
        $\vec{am} = (\vec{f})^{-1}(\vec{bn})$.
        D'où $(\vec{f^{-1}})_b (\vec{bn}) = (\vec{f})^{-1}(\vec{bn})$ d'où $(\vec{f^{-1}})_b = (\vec{f})^{-1}$.

        Or $f$ étant affine et bijective, $\vec{f}$ est linéaire et bijective et sa réciproque $(\vec{f})^{-1}$ est aussi linéaire.
        Donc $f^{-1}$ est affine d'application linéaire associée $(\vec{f})^{-1} = \vec{f^{-1}}$.
        \item Supposons que $f: \E \to \E$ avec $\dim(\E) < +\infty$
        Si $f$ est bijective, alors $f$ est injective et surjective.
        Réciproquement, si $f$ est injective alors $f$ est bijective car $f \text{ injective } \implies \vec{f} \text{ injective } \implies \vec{f} \text{ bijective } \implies f \text{ bijective }$.
    \end{enumerate}
\end{demonstration}

\begin{definition}
    Une application affine bijective est appelée un \textbf{isomorphisme affine}.
\end{definition}

\begin{demonstration}
    $f: \E \to \mathcal{F}$ une application affine et $\vec{f}: E \to F$ linéaire.

    Soit $\E' = a + E'$ un sous-espace affine de $\E$. Montrons que $f(\E')$ est un sous-espace affine de $\mathcal{F}$ passant par $f(a)$ et de direction $\vec{f}(E')$.

    C'est-à-dire $f(\E') = \{ m \in \mathcal{F} \mid \vec{f(a)m} \in \vec{f}(E')\}$

    Par double inclusion :
    Soit $m \in f(\E')$ alors $\exists m \in \E'$ tel que $m = f(m)$
    Or $m \in \E'$ donc $\vec{am} \in E'$
    d'où $\vec{f(a)m} \in \vec{f}(E')$ d'où l'inclusion $f(\E') \subset \{m \in \mathcal{F} \mid \vec{f(a)m} \in \vec{f}(E')\}$
    
    Réciproquement: soit $m \in \mathcal{F}$ tel que $\vec{f(a)m} \in \vec{f}(E')$
    Donc $\exists \vec{u} \in E'$ tel que $\vec{f(a)m} = \vec{f}(\vec{u})$
    Par la bijection $\funcdef{\theta_a}{\E}{E}{m}{\vec{am}}$,

    $\exists! m \in \E, \vec{u} = \vec{am}$ donc $\vec{f(a)m} = \vec{f}(\vec{am})$
    De plus $\vec{am} \in E' \implies m \in \E'$ donc $\vec{f(a)m} = \vec{f(a)f(m)} \implies m = f(m) \in f(\E')$
    \begin{hotwarn}
    À corriger
    \end{hotwarn}
\end{demonstration}

\end{document}